\chapter{Theoretical Foundations and Literature Review}
\label{ch:foundations}

This chapter sets out the five bodies of theory the thesis relies on and
reviews the empirical work that has tested them. \Cref{sec:discretetime}
covers discrete-time market design and the latency arbitrage problem it
addresses. \Cref{sec:adverse} covers information asymmetry and adverse
selection, which explain why a liquidity provider quotes a spread at
all. \Cref{sec:transparency} covers market transparency theory, which
governs the disclosure block of the design space.
\Cref{sec:auctiontheory} covers uniform-price auction theory, which
governs clearing and rationing. \Cref{sec:quality} covers the
measurement framework that turns all of this into numbers.\Cref{sec:priceimprovement} covers the price improvement effect and latency boundary arbitrage problem of Frequent Batch Auctions.


\section{Discrete-Time Market Design}
\label{sec:discretetime}

\subsection{The Mechanics of Latency Arbitrage }

The argument of \citet{budish2015hft} does not rest on any assumption
about irrationality or market power. It rests on a single structural
observation: in a continuous market, a public information event creates
a discrete arbitrage opportunity that is allocated by relative speed.

The information is continuously being discounted in the price which makes markets more efficient. Consider a liquidity provider quoting a two-sided market in an asset
whose fair value is tracked by a correlated, continuously observable
signal. When the signal jumps, every participant learns the same thing
at essentially the same instant. The liquidity provider wants to cancel
its stale quote; every other fast participant wants to hit it. Both
messages travel to the matching engine, and the engine processes them in
arrival order. Whether the stale quote is cancelled or picked off
depends on which message arrives first, which depends on network path
length, hardware, and code efficiency, not on anything informational.

Two consequences follow. First, the opportunity is mechanically
recurrent: it appears every time the public signal moves, which in a
liquid instrument is thousands of times per day. Second, the value of
being marginally faster than the next participant does not decline as
everybody invests in speed. Since only the winner is paid, the
equilibrium is a rent-dissipating race in which aggregate expenditure on
latency reduction is bounded by the aggregate value of the arbitrage
opportunities, and the opportunities themselves persist regardless of
how fast the market becomes. Speed is a purely relative good.

The cost is not borne by the racers. In equilibrium, a liquidity
provider that expects to lose a certain fraction of its volume to
picked-off quotes must recover that expected loss from the trades it
does not lose. This shifts the burden onto liquidity demanders in the
form of a wider spread, which \citet{budish2015hft} characterise as a
tax on liquidity.

\subsection{Empirical Magnitude}

\citet{aquilina2022quantifying} measure this directly. Using the London
Stock Exchange's message-level data, including messages that were never
executed, they identify latency races as clusters of near-simultaneous
attempts to take or cancel the same quote, and they observe both winners
and losers. Their findings are relevant here in three respects. Races
are extremely frequent and extremely short, so short that any batch
interval measured in milliseconds is orders of magnitude longer than the
time scale on which they are decided. A large share of the trading
volume that occurs at exactly the moment of a price change is race
volume. And the aggregate transfer, expressed as a share of trading
volume, is small in relative terms but large in absolute terms and
concentrated precisely in the effective spread paid by liquidity
demanders. The relevance for market design is that removing the race is
not merely an aesthetic improvement; it removes a component of the
spread.

\subsection{Batching as the Response}

Frequent batch auctions change two rules at once, and it is worth
keeping them analytically separate because they do different work.

\emph{Discrete time} removes the informational value of within-interval
arrival order. If the engine collects orders over $[t, t+\tau)$ and
matches them only at $t+\tau$, then any two orders inside the window are
symmetric with respect to time.

\emph{Uniform pricing} removes the value of queue position within the
matched set. If all executions in the batch occur at one price $\clr$,
then no participant obtains a better price by being earlier.

Either rule alone is insufficient. Discrete time with discriminatory
pricing would still reward earlier submission with a better price.
Uniform pricing in continuous time is not well defined. Together they
imply that speed pays only for crossing an interval boundary, that is,
for getting into batch $n$ rather than batch $n+1$.In frequent batch auctions the value of a
marginal nanosecond therefore falls from ``wins the race'' to ``matters
only in the vanishingly rare case that it changes which batch an order
lands in''. In the model of \citet{budish2015hft}, the equilibrium
consequence is that liquidity providers no longer need to price
sniping risk into their quotes and spreads narrow.

The residual cost of the mechanism is delay. If orders are distributed uniformly  an order takes on
average $\tau/2$ before execution, and the choice of $\tau$ is
accordingly a trade-off between the neutralisation of speed and the
delay imposed on everybody. This trade-off is one of the central design
questions of \Cref{ch:design}.

\section{Information Asymmetry and Adverse Selection}
\label{sec:adverse}

Latency arbitrage is one specific mechanism through which a liquidity
provider loses money to better-informed counterparties. The general
theory of that loss predates high-frequency trading by three decades and
is what makes the spread a meaningful economic object in the first
place.

\subsection{Glosten and Milgrom: The Spread as an Information Premium}

\citet{glosten1985bid} model a market maker who faces a mixed population
of informed and uninformed traders and who cannot tell them apart. The
market maker is competitive and risk-neutral, so it earns zero expected
profit. Because a buy order is more likely to come from an informed
trader when the asset is undervalued, every order carries information,
and the market maker must set its quotes so that they are conditionally
fair given the direction of the order it is about to receive. The result
is a bid--ask spread that exists purely as an information premium, even
with no inventory risk and no processing cost.

The comparative statics matter for this thesis. The spread widens with
the proportion of informed traders and with the precision of their
information, and it narrows when the market maker can identify a subset
of flow as uninformed. This is the theoretical foundation of order flow
segregation: if benign flow can be separated from potentially informed
flow, a liquidity provider can quote a tighter price against the benign
subset without cross-subsidising losses on the other.

\subsection{Kyle: Strategic Informed Trading and Market Depth}

\citet{kyle1985continuous} approaches the same problem from the side of
the informed trader. A single informed agent chooses how much to trade,
knowing that trading more reveals more. Noise traders provide the cover
that makes concealment possible, and a competitive market maker sets the
price as the conditional expectation of value given total order flow.
The equilibrium price function is linear, $p = \mu + \lambda \cdot y$,
where $y$ is aggregate order flow and $\lambda$ measures the price
impact per unit of order flow. Its reciprocal, market depth, is the
volume required to move the price by one unit.

Kyle's model is a batch model. The market maker in his setup observes
aggregate order flow and sets a single price, which is exactly the
structure of a uniform-price call auction, and the informed trader
spreads trading across periods to disguise it. The FBA literature can be
read as asking what happens when a real market is organised the way
Kyle's model already assumes. The parameter $\lambda$ also supplies the
empirical bridge: price impact is the operational measure of adverse
selection, and it is one of the metrics computed in
\Cref{sec:metrics}.

\subsection{Decomposing the Spread}
\label{sec:spread-decomposition}

The adverse selection premium that \citet{kyle1985continuous} and
\citet{glosten1985bid} place at the centre of the quoting decision is
not directly observable; only the total cost of a trade is. The standard strategy
is isolating it by the decomposition of the effective
spread into a realised spread and a price impact component
\citep{hasbrouck1993assessing,hendershott2011algorithmic}.

Let $p_k$ be the execution price of trade $k$, $D_k \in \{+1,-1\}$ its
direction (buyer- or seller-initiated), $m_k$ the midpoint immediately
before the trade, and $m_{k+\Delta}$ the midpoint a markout horizon
$\Delta$ afterwards. In proportional form,
\begin{align}
  S^{\mathrm{eff}}_k &= \frac{2 D_k (p_k - m_k)}{m_k},
    \label{eq:effspread}\\[3pt]
  S^{\mathrm{real}}_k(\Delta) &= \frac{2 D_k (p_k - m_{k+\Delta})}{m_k},
    \label{eq:realspread}\\[3pt]
  \mathrm{PI}_k(\Delta) &= \frac{2 D_k (m_{k+\Delta} - m_k)}{m_k},
    \label{eq:priceimpact}
\end{align}
so that $S^{\mathrm{eff}}_k = S^{\mathrm{real}}_k(\Delta) +
\mathrm{PI}_k(\Delta)$ trade by trade and in any volume-weighted
aggregate. The effective spread is what the liquidity demander paid
relative to the prevailing midpoint. The realised spread is what the
liquidity provider still holds once the market has revised its
valuation, and is therefore the gross compensation for supplying
liquidity. The price impact is the permanent revision of the midpoint in
the direction of the trade, the part of the spread the provider lost due to
the information the order carried.



In a public signal event like correlated prices (or other information to be discounted in the price), discrete uniform-price clearing removes
that mechanism, because within a batch the aggressive order and the
maker's cancellation are processed together and execution occurs at a
clearing price that already reflects the orders submitted in response to
the signal. It does not remove adverse selection in the theoretical
sense, since a genuinely informed trader can still price aggressively
into a batch \citep{eibelshauser2022frequent}; the prediction is a
reduction in the price impact component, not its elimination.



A fall in the effective spread driven by the price impact
component is the result the theory predicts. A fall driven by the
realised spread instead would mean liquidity providers simply earn less
for supplying liquidity, with the informational component untouched, and
would call for a different explanation. A fall in price impact with no
change in the effective spread would indicate that the benefit is
captured by providers rather than passed on, which is a statement about
competition among them rather than about the clearing rule.



\section{Market Transparency Theory}
\label{sec:transparency}

\citet{madhavan1996transparency} defines transparency as the ability of
participants to observe information about the trading process, and shows
that more of it is not unambiguously better. Greater pre-trade
transparency helps prices aggregate information and reduces uncertainty
about the state of the market, but it also exposes the trading intent of
participants who must transact in size, and it can increase price
volatility and reduce liquidity in some configurations. Consolidating
the argument with \citet{ohara1995microstructure}, disclosure design is
a genuine trade-off rather than a monotone improvement.

In an FBA the relevant question is narrower and sharper: what is
disclosed during the accumulation phase, before the auction clears? Four
regimes span the space.

\begin{itemize}
  \item \textbf{Fully open.} The complete set of accumulated orders is
  visible in real time. Price discovery is maximal, and so is exposure:
  a participant who submits early has broadcast its intent to everyone
  who will submit later in the same batch.
  \item \textbf{Indicative price only.} The engine publishes the
  clearing price that would result if the auction closed now, but not
  the underlying orders. This is the regime used by most equity opening
  and closing auctions.
  \item \textbf{Aggregate imbalance only.} The engine publishes the net
  excess demand or supply without revealing the price schedule.
  \item \textbf{Sealed bid.} Nothing is published until the auction
  clears.
\end{itemize}

The trade-off runs in a predictable direction. Moving toward opacity
protects submitted orders from being exploited but degrades the
information available for pricing, and classical auction theory predicts
that participants respond to uncertainty about the clearing price by
shading their bids \citep{klemperer1999auction}. Bid shading in a
uniform-price setting is not merely a transfer; it can reduce the volume
that clears, because shaded limit prices are less likely to cross.


This trade-off assumes the exchange operator can choose where on the previous options to
sit. Decentralized Blockchain venues supply the limiting case in which it cannot. A
transaction waiting in a public mempool is visible to everyone before it
executes, so such markets occupy the fully open regime by default,
without the benefit that regime is supposed to buy. The cost of
disclosure in \citet{madhavan1996transparency} is exposure of trading
intent, and in \citet{glosten1985bid} the informational disadvantage
falls on the quoter. A mempool inverts both: what is revealed is an
instruction that has not yet happened, and the party revealing it is the
one demanding liquidity. An observer can buy ahead of a pending swap and
sell into the price it moves, sandwiching the submitter. Batch clearing
removes this by construction rather than by rule, since if every order in
the interval settles at one uniform price, position within the batch
carries no value.


\section{Auction Theory and Uniform-Price Clearing}
\label{sec:auctiontheory}

\subsection{Why Uniform Pricing}

In a uniform-price auction, all winning participants transact at the same
market-clearing price, independent of the prices they submitted. A
participant's own bid therefore determines \emph{whether} it trades but,
except for the marginal order at which the aggregate bid and offer
schedules cross, not \emph{at what price}. Misreporting has no upside
under such a rule: bidding below one's valuation only risks losing a
trade that would have been acceptable, and bidding above it only risks
winning one that is not. For a price-taking participant with single-unit
demand, truthful revelation is therefore approximately optimal, a
property inherited from \citet{vickrey1961counterspeculation}. This is
the formal basis for the claim that batching substitutes competition in
price for competition in speed: if submitted limits are approximately
truthful, the clearing price aggregates valuations rather than reaction
times. Visibility of the accumulated batch does not disturb this, since
a participant who can anticipate the clearing price pays it whatever
limit it submits above it; what visibility changes is \emph{when} orders
are submitted, because a participant that can observe the indicative
price has reason to withhold its order until shortly before the batch
closes.

The qualification carried by ``approximately'' depends jointly on
pivotality and size, and separating the two cases clarifies what is at
stake. A participant too small to move the crossing point controls only
its own inclusion, so shading risks losing a trade and saves nothing;
truthful bidding is optimal without qualification. A pivotal participant
with single-unit demand does pin the clearing price with its own limit,
but the marginal unit is by definition the one carrying almost no
surplus, so the incentive to shade exists and is negligible. The case
that matters is the pivotal participant with multi-unit demand: its last
unit sets the price, and that price applies to every unit it buys, so
shading the marginal bid reduces the cost of the entire inframarginal
block while risking only the one unit that was barely worth trading. The
saving scales with size, the loss does not, and the incentive to
misreport is restored; \citet{klemperer1999auction} surveys this
literature.

Pivotality here is relative to the liquidity accumulated within one
batch rather than to the market as a whole, so the strength of the
truthfulness property is itself a function of the interval length.
Shorter intervals produce thinner batches, the same order that is atomistic
in a long window can be price-setting in a short one. 
This trade-off is discussed in \Cref{sec:interval}.

For every executed non-marginal order, uniform pricing also produces
price improvement relative to the submitted limit. A buyer willing to
pay $\pi_k$ who clears at $\clr < \pi_k$ captures $(\pi_k - \clr) q_k$
of surplus that, under continuous discriminatory execution at the
resting price, would in general have accrued differently. Aggregating
this quantity across executed orders gives the realised trader surplus
metric used in \Cref{sec:metrics}.

\subsection{The Clearing Problem}

The clearing price is not chosen by convention; it is the solution to an
optimisation problem. In the single-asset case, the standard rule
selects the price that maximises executed volume, with tie-breaking
rules for the case where an interval of prices achieves the same
maximum. The formal statement, including the multi-asset extension with
coherent cross-rates that
\citet{walther2021multitoken} and \citet{ramsay2023speedex} develop for
decentralised exchanges, is given in \Cref{sec:clearing}.

\subsection{Rationing and Induced Behaviour}

Maximising volume determines the price and the total quantity traded,
but when supply and demand are not exactly equal at $\clr$, the
mechanism must decide which orders on the long side are filled. This
is where auction theory has the most direct behavioural bite, and where
the market design literature is most explicit about unintended
consequences.

\citet{haynes2020precedence} study precedence rules in futures markets
and document that participants adapt their submission behaviour to the
allocation rule in measurable ways. The general pattern is that any rule
which makes allocation an increasing function of a variable the
participant controls creates an incentive to inflate that variable. If
allocation is proportional to submitted size, participants submit more
size than they want to trade. If allocation follows time priority within
the batch, participants compete to submit early, which reintroduces a
speed race at the scale of the interval. If allocation is random,
neither channel exists, but participants face execution uncertainty they
cannot manage. The design implications are developed in
\Cref{sec:rationing}.

\section{Measurement of Market Quality}
\label{sec:quality}

The empirical comparison relies on the standard microstructure toolkit.
Quoted spread and depth describe the state of the book and therefore the
cost a hypothetical trade would face. Effective spread describes what an
actual trade cost. Realised spread and price impact decompose that cost
into the part retained by the liquidity provider and the part lost to
adverse selection \citep{hasbrouck1993assessing}. \citet{amihud2002illiquidity}
provides a low-frequency summary of price impact per unit of volume that
is robust and easy to compare across mechanisms.
\citet{hasbrouck1995onesecurity} supplies the framework for attributing
price discovery when the same asset trades in several places, which is
relevant here because Hyperliquid publishes an oracle price that serves
as an external reference.

Two adaptations are required when these measures are applied to a batch
mechanism, and both are addressed explicitly in \Cref{sec:metrics}.
First, several measures reference a prevailing midpoint at the moment of
the trade, but in an FBA all trades in an interval occur at one instant,
so the reference must be defined as the midpoint at the start of the
batch. Second, the quoted spread has no direct analogue during
accumulation, and is defined here as the implied spread between the
marginal unexecuted buy and the marginal unexecuted sell of the batch.

Finally, the experimental convention this thesis follows is the one
established by \citet{gode1993allocative}: to compare institutions
meaningfully, hold the induced order flow constant and vary only the
mechanism. Their demonstration that market institutions can produce
allocative efficiency largely independently of trader rationality is
also a methodological licence for replay-based simulation, since it
implies that a substantial part of what a mechanism does is mechanical
rather than behavioural. \citet{lebaron2006agent} surveys the
agent-based alternative, in which trader behaviour is modelled and
allowed to adapt; \Cref{sec:limitations} discusses why replay is the
appropriate choice here and what it gives up.

\section{Price Improvement as a Distributive Artifact and the Boundary
Repricing Race}
\label{sec:priceimprovement}

The argument developed in this section is, to the author's knowledge, not
stated in this form in the existing literature. Its components are
individually familiar: the residual value of low latency near a batch
boundary is established by \citet{schlegel2023transaction}, the tendency
of sophisticated traders to submit late in transparent call auctions is
documented by \citet{daures2026click}, and the strategic consequences of
the clearing-price selection rule in uniform-price auctions are studied
by \citet{burkett2020uniform}. What has not been made is the connection
between them. This section argues that the price improvement conferred
by uniform clearing, which the frequent batch auction literature reports
as a benefit of the design, is in fact a transfer whose size is set by a
manipulable rule, and that competition over that transfer is itself
sufficient to generate a timing race at the batch boundary, independently
of the informational motives usually invoked to explain late submission.


The literature on frequent batch auctions treats the price improvement
that uniform clearing confers on inframarginal orders as a benefit of
the design. This section argues that the treatment rests on a category
error. Price improvement measured on one side of the book is not a
welfare quantity but a transfer, and the rule that determines its size
is manipulable by the timing of order submission. The consequence is
that the FBA does not remove the incentive to compete on speed; it
converts a race for priority into a race for proximity to the batch
boundary. The argument is developed formally below and then bounded:
it is a qualification of the Budish, Cramton and Shim mechanism, not a
refutation of it.

\subsection{Setup}

Consider a single batch collecting orders over the window $[0,\tau]$.
Let $\mathcal{B}$ and $\mathcal{S}$ denote the sets of buy and sell
orders submitted in the window, each order $k$ characterised by a limit
price $\pi_k$, a quantity $q_k$, and a submission time $t_k$. Define the
aggregate schedules
\begin{equation}
  D(p) = \sum_{k \in \mathcal{B}} q_k \mathbf{1}\{\pi_k \ge p\},
  \qquad
  S(p) = \sum_{k \in \mathcal{S}} q_k \mathbf{1}\{\pi_k \le p\},
\end{equation}
and the executable volume $V(p) = \min\{D(p), S(p)\}$. The clearing rule
selects a price from the volume-maximising set
\begin{equation}
  \mathcal{P}^{*} = \arg\max_{p} V(p),
  \qquad V^{*} = \max_p V(p),
\end{equation}
which, since $D$ is non-increasing and $S$ non-decreasing, is a closed
interval $[\underline{p}, \overline{p}]$ whenever the book crosses. When
this interval is non-degenerate the clearing rule is not fully
determined by the volume objective, and a secondary selection rule
$\sigma : \mathcal{P}^{*} \to \mathbb{R}$ must be specified. Real
auctions resolve the indeterminacy by minimising the residual imbalance
and then by proximity to a reference price
\citep{haynes2020precedence}. Let $x_k \ge 0$ denote the executed
quantity of order $k$ under a volume-maximising allocation, so that
$\sum_{\mathcal{B}} x_k = \sum_{\mathcal{S}} x_k = V^{*}$.

\subsection{Price Improvement Is a Transfer}

Write the realised surplus of the two sides at clearing price $p$ as
\begin{equation}
  \mathrm{PI}^{\mathcal{B}}(p)
    = \sum_{k \in \mathcal{B}} (\pi_k - p)\, x_k,
  \qquad
  \mathrm{PI}^{\mathcal{S}}(p)
    = \sum_{k \in \mathcal{S}} (p - \pi_k)\, x_k .
\end{equation}

\begin{lemma}[Invariance of total surplus]
\label{lem:invariance}
For any $p \in \mathcal{P}^{*}$ with allocation $x$ held fixed,
\begin{equation}
  W = \mathrm{PI}^{\mathcal{B}}(p) + \mathrm{PI}^{\mathcal{S}}(p)
    = \sum_{k \in \mathcal{B}} \pi_k x_k
      - \sum_{k \in \mathcal{S}} \pi_k x_k ,
\end{equation}
which does not depend on $p$.
\end{lemma}


Expanding the two sums, the terms in $p$ enter as
$p\left(\sum_{\mathcal{S}} x_k - \sum_{\mathcal{B}} x_k\right)$, which
vanishes by flow conservation.


Two consequences follow immediately. First,
$\partial \mathrm{PI}^{\mathcal{B}} / \partial p = -V^{*}$ and
$\partial \mathrm{PI}^{\mathcal{S}} / \partial p = +V^{*}$: every unit of price improvement conferred on
one side is subtracted from the other. Second, and consequently, a
measured increase in buy-side price improvement is not evidence that the
mechanism created value. Total gains from trade are fixed by the
allocation, and the clearing price only divides them. The metric
reported in \Cref{sec:metrics} must therefore be read as a distributive
statistic, and reported for both sides jointly rather than for one.

This reframing is what motivates the rest of the section. If
$\mathcal{P}^{*}$ were a singleton, the division would be pinned down by
the submitted limits and there would be nothing to contest. Because it
is generically an interval, the surplus split depends on a rule, and any
participant able to influence the position of $\mathcal{P}^{*}$ or of
$\sigma$ within it has an incentive to do so.

\subsection{The Last-Mover Advantage}

Suppose the engine operates under the fully open disclosure regime of
\Cref{sec:transparency}, so that a participant submitting at time $t$
observes the schedules formed by all orders with $t_k < t$.

\begin{proposition}[Late submission weakly dominates]
\label{prop:lastmover}
Let a participant observe the residual schedule it faces and let
$\sigma$ be weakly monotone in submitted limit prices. Then the
participant's optimal limit is the one that moves $\sigma(\mathcal{P}^{*})$
to the boundary of $\mathcal{P}^{*}$ most favourable to it subject to
remaining executed, and its realised surplus is weakly increasing in the
information available at submission. Submitting later therefore weakly
dominates submitting earlier.
\end{proposition}

The argument is the standard one for sequential-move bargaining over a
fixed surplus, applied here to the division established in
\Cref{lem:invariance}. A seller who observes a resting buy limit of
$\pi_{\mathcal{B}}$ knows that any ask weakly below it will execute, and
therefore submits just below $\pi_{\mathcal{B}}$ rather than at its own
valuation, capturing the interval. The truthfulness property of
\Cref{sec:auctiontheory} is not violated in its own terms, since it is
stated for a participant that cannot move the clearing price; what
disclosure does is make the boundaries of $\mathcal{P}^{*}$ computable,
and hence make deviation actionable for any participant whose order is
large enough to reach them.

Two conditions therefore expose a participant to being the donating side
of the transfer: submitting without observing the accumulated batch, as
an uninformed participant does, and submitting early without
subsequently repricing, which an informed participant may do by choice
or by failure to cancel in time. In either case the resting limit is
available to be cleared against at a price determined by orders that
arrive afterwards. Since the transfer is zero-sum by
\Cref{lem:invariance}, every participant exposed in this way is matched
by one with an incentive to identify and exploit that exposure.



An order resting in an open batch is thus a free option written to every
subsequent submitter in the same window, exercisable at the submitter's
own limit price and, absent a cancellation, not repriceable. This is the
\citet{copeland1983information} characterisation of a limit order,
transposed from the continuous book, where the option is written to
traders reacting to news, into the batch, where it is written to traders
reacting to the accumulating batch itself.

\subsection{Cancellation, Repricing, and the Deadline Race}

The natural defence against \Cref{prop:lastmover} is to withdraw or
reprice an order that has become adversely positioned. This is where the
speed race re-enters.

\begin{proposition}[Latency monotonicity]
\label{prop:latency}
Let participant $i$ face message latency $\ell_i > 0$ and let the batch
close deterministically at $\tau$. A message dispatched at time $t$ is
included in the batch if and only if $t + \ell_i \le \tau$, so the
latest observable state on which $i$ can act is that prevailing at
$\tau - \ell_i$. The information set available to $i$ at the moment of
its final decision is therefore ordered by $\ell_i$, and expected
surplus is non-increasing in $\ell_i$.
\end{proposition}

The structure of the advantage differs from the continuous case but its
source does not. In a continuous book, lower latency wins a contest for
priority: the fast trader is served first. In a batch, priority within
the window is irrelevant by construction, and lower latency instead buys
a later decision point. The fast trader observes more of the batch
before committing, can reprice on more recent information, and can place
its final message closer to the boundary while retaining the same
probability of inclusion. Speed is not neutralised; its payoff is
relocated from ordering to deadline precision. That some speed advantage
survives batching is consistent with the experimental evidence of
\citet{aldrich2020experiments}, who find that investment in speed falls
but does not disappear under batch clearing, and with the model of
\citet{zoican2021speed}.

Three remarks bound the mechanism. First, prohibiting cancellation
within the batch does not close the channel, because a participant with
two-sided access can neutralise an adverse position by submitting an
offsetting order, which functions as a synthetic cancellation while also
perturbing $D$ and $S$; enforcement would require self-match prevention
at the account level. Second, the prohibition would in any case
intensify rather than relieve the problem, since it raises the cost of
early submission and therefore strengthens the incentive identified in
\Cref{prop:lastmover}. Third, the channel is disclosure-conditional: the
informational component of the last-mover advantage requires that the
accumulating batch be observable, and it is absent under a sealed-bid
regime. What survives sealed bidding is only the incentive to submit
late in order to condition on external information, which is the
boundary-sniping problem already identified by
\citet{budish2015high}.

\subsection{Equilibrium Implication and Design Responses}

If late submission weakly dominates and lateness is bounded only by
$\ell_i$, submissions concentrate near the boundary and are ordered by
latency. In the limit the accumulation phase carries no orders and
therefore no information, and the auction degenerates into a sealed
tender conducted in the final microseconds of the window: the batch
retains its length nominally while losing the property the length was
chosen to deliver.

Three design responses follow directly, and each is a parameter of the
design space developed in \Cref{ch:design}. Randomising the close time
makes waiting costly, since an order dispatched at $t$ is included only
if the realised $\tau$ exceeds $t + \ell_i$; the participant then solves
an interior optimum trading the value of information against the risk of
exclusion, which supplies a justification for randomisation independent
of the sniping argument usually advanced for it. Restricting disclosure
to an aggregate imbalance or to nothing at all removes the informational
component of \Cref{prop:lastmover} while leaving bid shading, so the
choice is between the two costs rather than between one cost and none.
Finally, and least discussed in the literature, the selection rule
$\sigma$ can be made insensitive to the marginal submitted limits, for
instance by choosing the point of $\mathcal{P}^{*}$ closest to an
external reference price rather than a function of the interval
endpoints. This removes the lever without altering the allocation, since
by \Cref{lem:invariance} the allocation is invariant on $\mathcal{P}^{*}$
in any case.



